% Generated draft: requires mathematical and source review.
\documentclass[11pt]{article}
\usepackage[margin=25mm]{geometry}
\usepackage{fontspec}
\setmainfont{Latin Modern Roman}
\usepackage{amsmath,amssymb,mathtools,unicode-math}
\setmathfont{Latin Modern Math}
\usepackage{enumitem,needspace,xurl,hyperref}
\setlength{\parindent}{0pt}
\setlength{\parskip}{5pt}
\setlength{\emergencystretch}{3em}
\newcommand{\N}{\mathbb{N}}
\newcommand{\Z}{\mathbb{Z}}
\newcommand{\Q}{\mathbb{Q}}
\newcommand{\R}{\mathbb{R}}
\newcommand{\C}{\mathbb{C}}
\newcommand{\F}{\mathbb{F}}
\newcommand{\A}{\mathbb{A}}
\newcommand{\PP}{\mathbb{P}}
\newcommand{\E}{\mathbb{E}}
\newcommand{\eps}{\varepsilon}
\newcommand{\dd}{\,\mathrm{d}}
\begin{document}
\section*{UnsolvedMath ID: 2306054}
\textbf{Source problem code:} catalogue entry
\subsection{Definitions and mathematical statement}
Let $D\subset\C$ be a Jordan domain with boundary $C$. Under the usual bounded-interior convention, let $\Phi$ conformally map $\widehat\C\setminus\overline D$ onto $\{|w|>1\}$, normalized at infinity with positive leading coefficient. A standard normalization of the Faber polynomial $F_n$ is the polynomial part of the Laurent expansion of $\Phi(z)^n$ at infinity. The source does not state this normalization explicitly.

Using the source's displayed series, let $S(D)$ consist of univalent functions on $D$ admitting
\[
f(z)=F_1(z)+\sum_{n=2}^{\infty}a_nF_n(z).
\]
Here univalent means holomorphic and injective. Convergence and uniqueness of the Faber representation must be specified; local uniform convergence is a natural analytic interpretation, not an extra assertion that every univalent function on every Jordan domain has this representation. No constant $F_0$ term or free $F_1$ coefficient is added.

For each $n\ge2$, its coefficient region is
\[
\mathcal R_n(D)=\{a_n\in\C:f\in S(D)\}.
\]
Lucas asks whether at least one coefficient region has the ``same shape'' as $C$, or whether at least a star-like assertion holds. The existence quantifier is over a coefficient index, potentially depending on $D$, rather than a claim about every coefficient. The record does not define same shape: comparison of a region with the boundary curve $C$ might concern its boundary up to similarity or another specified equivalence, but no such equivalence is supplied. Nor is the star-like clause well typed as written: it may ask whether a coefficient region is star-shaped or ask for a statement after restricting $S(D)$ to star-like functions. A set $K$ is star-shaped about $p$ if every segment joining $p$ to a point of $K$ is contained in $K$; a star-like univalent function usually has image star-shaped about a specified point. The source fixes neither reading nor center.

For context, the source attributes to Royster a statement for a star-like $f$ in an ellipse: there are an angle $\theta_f$ and sequences $\lambda_n>\mu_n>0$ for $n\ge2$ such that $a_n$ lies in an ellipse centered at zero, with major and minor axes $\lambda_n,\mu_n$, inclined at $\theta_f$. The same angle is used for all coefficients of that function. The numerical axes are said to be known but are not supplied, and the source does not fix full-axis versus semiaxis conventions or the allowed dependence of these quantities. This containment is not asserted to identify the exact coefficient region of the unrestricted class.

The supplied triage says the question was open in the 2018 edition and no definitive resolution was located. This is an attributed source assessment, not independently verified openness.
\subsection{Short English statement}
For univalent functions expanded in a domain's Faber polynomials, does at least one coefficient region reflect the boundary's shape, or satisfy an appropriate star-like property? The geometric equivalence and the star-like clause are undefined. The reported ellipse bounds concern a restricted star-like class and individual coefficients, so they do not supply an exact coefficient-region answer.
\subsection{Sources}
\begin{itemize}
\item Local source registry: \url{data/unsolvedmath/problems.json}, permanent UnsolvedMath ID $2306054$, source code catalogue entry, including the supplied OPEN-TRIAGE assessment. Dataset attribution: Ulam AI and UnsolvedMath Contributors; original-source attribution and rights are retained.
\item Hayman, \emph{Research Problems in Function Theory}, source-list date 2018, Problem 6.54, attributed to K. W. Lucas: \url{https://arxiv.org/abs/1809.07200}. The metadata reports source-TeX extraction and July 29, 2026 access, with current-status and rights review flags.
\item Royster's ellipse coefficient-containment result is quoted in the statement, but no full citation, formulas for the known axes, or exact star-like convention is supplied. No linked source was inspected for this draft.
\end{itemize}
\end{document}
